% 第 1 课时：集合的含义与表示
\chapter{集合的含义与表示}

\begin{objective}
\begin{itemize}
  \item 了解集合的含义，体会元素与集合之间的“属于”关系；
  \item 掌握列举法与描述法，并能根据问题选择恰当的表示方法；
  \item 理解集合中元素的确定性、互异性与无序性。
\end{itemize}
\end{objective}

\begin{keypoint}
重点：集合的表示方法，元素与集合的关系。\\
难点：描述法中代表元素与约束条件的辨认，例如 $\{x\in\mathbb{R}\mid x^2-1=0\}$ 与 $\{x^2-1=0\}$ 的区别。
\end{keypoint}

\begin{procedure}
\section{情境引入}
考察下列对象：
\begin{enumerate}
  \item 我校高一年级全体学生；
  \item 方程 $x^2-1=0$ 在实数范围内的解；
  \item 所有大于 $3$ 的整数。
\end{enumerate}
它们的共同特征是：都是一些\emph{确定的、互不相同的}对象的整体。由此给出集合的概念。

\section{概念形成}
\begin{definition}
一般地，我们把研究对象统称为\emph{元素}，把一些元素组成的总体叫做\emph{集合}（简称\emph{集}）。
\end{definition}
\begin{itemize}
  \item 元素与集合的关系：$a\in A$ 或 $a\notin A$；
  \item 集合相等：构成两个集合的元素完全一样，则 $A=B$；
  \item 常用数集：$\mathbb{N}$、$\mathbb{N}^*$、$\mathbb{Z}$、$\mathbb{Q}$、$\mathbb{R}$。
\end{itemize}

\section{例题辨析}
\begin{example}
用列举法表示集合 $\{x\in\mathbb{Z}\mid -2<x\le 3\}$。
\end{example}
\begin{solution}
逐一检验整数：$-1,0,1,2,3$，故为 $\{-1,0,1,2,3\}$。
\end{solution}

\begin{example}
集合 $A=\{x\mid x^2=a\}$ 中元素的可能情况取决于 $a$ 的取值，试讨论。
\end{example}
\begin{solution}
当 $a<0$ 时 $A=\varnothing$；当 $a=0$ 时 $A=\{0\}$；当 $a>0$ 时 $A=\{-\sqrt{a},\sqrt{a}\}$。
\end{solution}

\section{课堂练习}
\begin{exercise}
用适当的方法表示下列集合：
\begin{enumerate}
  \item 由 $1,2,3,4,5$ 组成的集合；
  \item 函数 $y=x^2-1$ 的图象与 $x$ 轴交点的横坐标。
\end{enumerate}
\end{exercise}
\end{procedure}

\begin{board}
\begin{itemize}
  \item 集合三要素：确定性、互异性、无序性；
  \item 表示方法：列举法 $\{1,2,3\}$，描述法 $\{x\mid p(x)\}$；
  \item 代表元素决定集合含义。
\end{itemize}
\end{board}

\begin{homework}
\begin{enumerate}
  % TODO(待核对教材)：原写「教材习题 1.1 第 $1\sim4$ 题」，题号为占位，暂隐去
  \item 复习本节内容，完成课后练习；
  \item 思考：$\{1,1,2\}$ 能否表示一个集合？为什么？
\end{enumerate}
\reading{拓展阅读 01}{集合之美}
\end{homework}
